\intobmk\chapter*{术语中英对照}

\begingroup
\renewcommand{\nomenclatureitem}[3][ ]{%
    \noindent
    \parbox[t]{0.23\textwidth}{#2}\hfill
    \parbox[t]{0.73\textwidth}{#3\hfill#1}\par
    \vspace{0.4ex}%
}

\nomenclatureitem{\textbf{英文术语或专名}}{\textbf{中文名称与本文含义}}

\nomenclatureitem{AI}{Artificial Intelligence，人工智能；辅助检索、归约构造、实例分析与证明整理的方法}
\nomenclatureitem{LLM}{Large Language Model，大语言模型；根据问题定义和交互信息生成文本或程序的模型}
\nomenclatureitem{Code Agent}{代码智能体；可编辑文件、执行程序并依据结果继续交互的智能体}
\nomenclatureitem{Agent Harness}{智能体运行与评测框架；组织任务、工具、运行状态和反馈的外部环境}
\nomenclatureitem{Oracle}{实例判定器；在评价端返回问题实例 YES/NO 结果的判定程序}
\nomenclatureitem{Skill}{技能模块；向智能体提供任务知识和操作规范的说明模块}
\nomenclatureitem{ReAct}{推理与行动协同方法；使模型交替执行推理、操作和观察的方法}
\nomenclatureitem{SWE-bench}{软件工程任务基准；以真实代码库问题评价程序修复能力的基准}
\nomenclatureitem{SWE-agent}{软件工程智能体；面向代码库浏览、编辑和测试执行的智能体系统}
\nomenclatureitem{OpenHands}{软件开发智能体平台；提供终端、浏览器和受限运行环境的系统}
\nomenclatureitem{DeepSeek}{DeepSeek-V4-pro，早期消融实验使用的大语言模型}

\endgroup
